\section{Conclusion}
\label{sec:conclusion}

Multiple instance learning is the masked-data identifiability
problem from reliability statistics, with instances playing the
role of components, bags playing the role of candidate sets, the
``at least one positive instance'' rule playing the role of series
(logical-OR) aggregation, and singleton bags playing the role of
singleton candidate sets. We proved (i) a rank condition on the
bag-by-instance-type composition matrix that determines when
instance-type positivity is identifiable from bag labels, (ii) a
bag-prevalence consistency theorem explaining why a model that
predicts bag labels well can still mis-rank instances, and (iii) a
first-order bias bound under aggregation-rule misspecification that
recovers the standard-versus-collective assumption taxonomy as a
classification of how aggregation rules depart from the single-cause
OR structure, with threshold and proportion rules preserving the C1
coarsening condition and emergent rules violating it.

The framework places Diverse Density, mi-SVM, attention-based MIL,
and CLAM in a common assumption-naming language and supplies a
precise vocabulary for when each method's coarsening assumptions
hold. It does not by itself derive new estimators; the contribution
is diagnostic and structural rather than algorithmic. A confound between
intrinsic instance positivity and the noisy-OR firing rate is the
discrete-label analogue of the spike-in capture-efficiency gap in
scRNA-seq, and it shows that absolute instance calibration requires
instance-resolved auxiliary labels.

The principal open extensions are continuous instance features,
within-bag instance dependence, and identifiability of deep-pooling
architectures whose practical behaviour may exceed the linear rank
condition. Each is a natural follow-up, and each fits the same
coarsening scaffold shared with the scRNA-seq and
spatial-transcriptomics applications \citep{towell2026scrnacoarsening,
towell2026spatialcoarsening}.
